% ============================================================
% Section 186: 各种电子组态的光谱项与简并度
% ============================================================
\section{量子力学：各种电子组态的光谱项与简并度}
\label{sec:atomic-configurations-terms}

\begin{modelbox}[题目：各种电子组态的光谱项与简并度]
计算下列电子组态的 $^{2S+1}L_J$ 值，并证明每个光谱项的总态数等于组态简并度。
\begin{enumerate}
    \item[(1)] $2s2p$；\quad (2) $2p3p$；\quad (3) $(2p)^2$；\quad (4) $(3d)^{10}$；\quad (5) $(3d)^9$。
\end{enumerate}
\end{modelbox}

\begin{tipbox}[考点快速分析]
\begin{itemize}
    \item \textbf{非同科电子}：所有 $L$、$S$ 组合均允许，直积后反对称化
    \item \textbf{同科电子}：Pauli原理限制空间与自旋对称性相反
    \item \textbf{满壳层}：$L=S=J=0$，对总角动量无贡献
    \item \textbf{空穴等效}：$(nl)^{4l+2-k}$ 与 $(nl)^k$ 的光谱项相同
    \item \textbf{态数验证}：$\sum(2J+1)=$ 组态简并度
\end{itemize}
\end{tipbox}

\begin{derivationbox}[标准解题过程]
\textbf{(1) $2s2p$（非同科）}

$l_1=0,l_2=1\Rightarrow L=1$；$s_1=s_2=1/2\Rightarrow S=1,0$。

光谱项：$^3P_{2,1,0}$、$^1P_1$。

态数：$(5+3+1)+3=12$，等于 $2\times6=12$（$2s$ 有 2 态，$2p$ 有 6 态）。

\textbf{(2) $2p3p$（非同科）}

$l_1=l_2=1\Rightarrow L=2,1,0$；$S=1,0$。所有组合均允许。

光谱项：$^3D_{3,2,1}$、$^3P_{2,1,0}$、$^3S_1$、$^1D_2$、$^1P_1$、$^1S_0$。

态数：$15+9+3+5+3+1=36$，等于 $6\times6=36$。

\textbf{(3) $(2p)^2$（同科）}

总简并度 $\binom{6}{2}=15$。两 $p$ 电子轨道耦合 $L=2,1,0$，空间对称性：$L=2$ 对称、$L=1$ 反对称、$L=0$ 对称。

同科电子要求空间与自旋对称性相反：
\begin{itemize}
    \item $L=2$（对称）$\leftrightarrow S=0$：$^1D_2$（5 态）
    \item $L=1$（反对称）$\leftrightarrow S=1$：$^3P_{2,1,0}$（$5+3+1=9$ 态）
    \item $L=0$（对称）$\leftrightarrow S=0$：$^1S_0$（1 态）
\end{itemize}

总态数 $5+9+1=15$，一致。光谱项：$\boxed{^1D_2,\;^3P_{2,1,0},\;^1S_0}$。

\textbf{(4) $(3d)^{10}$（满壳层）}

满壳层唯一占据方式，简并度 1。所有角动量配对，$L=S=J=0$。

光谱项：$\boxed{^1S_0}$，态数 $2J+1=1$。

\textbf{(5) $(3d)^9$（满壳层缺一个电子）}

简并度等于空穴态数 $=10$。满壳层 $L=S=0$，去掉一个 $d$ 电子等价于加上一个 $d$ 空穴，$L=2,S=1/2$。

$J=5/2,3/2$。光谱项：$\boxed{^2D_{5/2},\;^2D_{3/2}}$。

态数：$6+4=10$，一致。
\end{derivationbox}

\begin{formulabox}[答案汇总]
\begin{center}
\begin{tabular}{llll}
\toprule
\textbf{组态} & \textbf{简并度} & \textbf{光谱项} & \textbf{态数和} \\
\midrule
$2s2p$ & 12 & $^3P_{2,1,0},\;^1P_1$ & 12 \\
$2p3p$ & 36 & $^3D,\;^3P,\;^3S,\;^1D,\;^1P,\;^1S$ & 36 \\
$(2p)^2$ & 15 & $^1D_2,\;^3P_{2,1,0},\;^1S_0$ & 15 \\
$(3d)^{10}$ & 1 & $^1S_0$ & 1 \\
$(3d)^9$ & 10 & $^2D_{5/2},\;^2D_{3/2}$ & 10 \\
\bottomrule
\end{tabular}
\end{center}
\end{formulabox}

\begin{insightbox}[物理图像]
\begin{itemize}
    \item 同科电子的光谱项推导是原子物理的核心技能。$(2p)^2$ 的 $^1D,\;^3P,\;^1S$ 是Hund定则的经典例证：$S$ 最大（$^3P$）的态能量最低，因为自旋平行使空间波函数反对称，减小库仑排斥。
    \item 满壳层规则极大简化了重原子的角动量耦合计算：只需考虑价电子或空穴。
    \item 空穴等效原理源于满壳层的$SO(3)$标量性质。在固体物理中，该原理推广为``粒子-空穴对称''，是理解半导体和金属电子结构的基础。
\end{itemize}
\end{insightbox}
